\subsection{Insights and Design Implications}
\label{sec:design-insights}

Three premises motivate the design (Figure~\ref{fig:design-insights}).
The first draws on published agent-prediction results; the other two
characterize selected sparse-attention releases and recovery demand.

\begin{figure*}[t]
  \centering
  \captionsetup[subfigure]{font=normalsize,labelfont=normalfont,
    textfont=normalfont,justification=centering,singlelinecheck=false,skip=3pt}
  \begin{subfigure}[t]{0.32\textwidth}
    \vspace{0pt}
    \includegraphics[width=\linewidth]{figures/figure3a-agent-prediction.pdf}
    \caption{Agent trajectory prediction}
    \label{fig:insight-agent-time}
  \end{subfigure}\hfill
  \begin{subfigure}[t]{0.32\textwidth}
    \vspace{0pt}
    \includegraphics[width=\linewidth]{figures/figure3b-sparse-models.pdf}
    \caption{Selected sparse releases}
    \label{fig:insight-sparse-adoption}
  \end{subfigure}\hfill
  \begin{subfigure}[t]{0.32\textwidth}
    \vspace{0pt}
    \includegraphics[width=\linewidth]{figures/figure3c-retry-kv-prediction.pdf}
    \caption{Retry KV prediction}
    \label{fig:insight-retry-kv}
  \end{subfigure}
  \caption{Design evidence. (a) Adapted from Heddle~\cite{zhang2026heddle},
    Figure 13: Pearson correlation of predicted and actual remaining
    trajectory length after one and two agent steps, with Qwen3-14B.
    Values are extracted from the published vector figure; these are
    literature results, not AegisKV runtime-error measurements.
    (b) Selected releases and
    native context: 1 MiniCPM4, 2 MiniCPM4.1, 3 DeepSeek-V3.2-Exp,
    4 DeepSeek-V3.2, 5 GLM-5, 6 MiniCPM-SALA, 7 DeepSeek-V4 preview,
    8 MiniMax-M3, 9 GLM-5.2, 10 LongCat-Flash-Lite-Sparse.
    Dates follow announcements/model cards. Gray steps trace the sample
    maximum; markers denote top-$k$ units in sparse branches.
    (c) Three native PG19 captures/model: 32K history; queries
    32,257--32,768 predict distinct historical KV for 32,769--32,832.
    Objects are layer/64-token blocks containing all KV heads; both
    policies include mandatory blocks in the 40\% budget.
    Seer uses an experimental Qwen3-8B Prefill gate~\cite{song2026seerqwen3}
    (threshold $10^{-3}$; final-block dense override disabled).
    Bars/whiskers are document means/ranges, not fault-recovery runs.}
  \label{fig:design-insights}
  \Description{Three side-by-side panels show Heddle's published agent
    trajectory prediction correlations, a primary-source survey of sparse model releases,
    and held-out retry KV coverage. Past-frequency ranking improves
    coverage substantially for NOSA and only modestly for MiniCPM4
    and SeerAttention-Qwen3.}
\end{figure*}

\paragraph{Insight 1: execution context predicts remaining agent work.}
Existing predictors exploit task semantics and execution progress.
Heddle~\cite{zhang2026heddle} reports trajectory-length correlations of
approximately 0.72, 0.88, and 0.94 for coding, search, and math agents
after two steps, improving over one-step predictions
(Figure~\ref{fig:design-insights}(\subref{fig:insight-agent-time})).
AgentDecarbonizer~\cite{yan2026agentdecarbonizer} additionally estimates
task durations and updates remaining time from observed progress.
AegisKV builds on these capabilities: the Tracker converts remaining-time
estimates into workflow slack for KV protection. Trajectory correlation
supports predictability; runtime calibration remains deployment-specific.

\paragraph{Insight 2: sparse attention changes the recovery objects.}
Long-context models use sparse mechanisms including InfLLMv2,
DSA, or compressed/hybrid attention~\cite{minicpm4model,
openbmb2025minicpm41,deepseek2025v32exp,deepseek2025v32model,
zai2026glm5model,openbmb2026sala,deepseek2026v4model,minimax2026m3,
zai2026glm52,longcat2026sparse}.
Figure~\ref{fig:design-insights}(\subref{fig:insight-sparse-adoption})
establishes release examples, not deployment share.
The KV Master therefore consumes actual layer/block selections and
model-specific recovery state instead of valuing all KV bytes equally.

\paragraph{Insight 3: past selections predict part of retry demand.}
With 40\% of each model's historical KV bytes, historical frequency
covers 77.8\%, 53.6\%, and 43.0\% of distinct needed KV for NOSA-3B,
MiniCPM4-8B, and SeerAttention-Qwen3-8B, versus 59.7\%, 50.9\%, and
42.0\% for recency
(Figure~\ref{fig:design-insights}(\subref{fig:insight-retry-kv})).
Seer's demand spans 92.5\% of historical blocks after the head union,
limiting even oracle placement to 43.3\% coverage. Thus gains depend on
the model and placement granularity. These held-out book-text suffixes
motivate advance placement with exact demand fetches for misses;
they do not establish performance on agentic failure retries.
